\documentclass[12pt]{article}
\usepackage[utf8]{inputenc}
\usepackage[T1]{fontenc}
\usepackage{hyperref}
\title{Tutelo Language Resources for K-12 Curriculum Development}
\author{Compiled by Claude Sonnet 5 (claude-sonnet-5) for Andr\'as Kornai}
\date{}
\begin{document}
\maketitle

\section{Introduction}
Tutelo (ISO 639-3 tta), also called Tutelo-Saponi, is a Siouan language once spoken in Virginia and West Virginia. It has no native speakers. Sources date its end to the death of Nikonha in 1871 or of Albert Green in the 1980s. Horatio Hale, J. O. Dorsey, J. N. B. Hewitt, Edward Sapir and Marianne Mithun recorded it. A revival effort is reported from around 2019.

\section{Primary Educational Resources}
No school curricula or primers were found in search results.

\section{Linguistic and Reference Materials}
\begin{itemize}
\item Giulia Oliverio (1997), \emph{A Grammar and Dictionary of Tutelo}, University of Kansas dissertation: \url{https://kuscholarworks.ku.edu/entities/publication/46f6622c-4d2c-4647-a4a5-8781193a72ba}
\item Glottolog record for Oliverio (1997): \url{https://glottolog.org/resource/reference/id/40571}
\item Corey Roberts (2020), a sketch grammar of Tutelo-Saponi for revitalization, University of Arizona: \url{https://repository.arizona.edu/handle/10150/648661}
\item Wikipedia article on the Tutelo language: \url{https://en.wikipedia.org/wiki/Tutelo_language}
\end{itemize}

\section{Religious and Community Resources}
None found in search results.

\section{Digital and Audio Resources}
\begin{itemize}
\item Tutelo-Saponi (Monacan) Living Dictionary. Its About page says it was started in 2021 with Endangered Language Fund support and led by Dr. Marvin Richardson, for Monacan, Haliwa-Saponi, Occaneechi, Sappony and Ohio Saponi communities: \url{https://livingdictionaries.app/tutelo-saponi}
\end{itemize}

\section{Curriculum Development Context}
The language is known only from historical records, so any curriculum would be built on these. It should be done with the descendant communities, using the Living Dictionary and the two grammars as a base.

\end{document}